\begin{itemize}
  \item What is your chosen human rights issue and how did it begin? \\
      \bf{I chose Human Trafficking in Thailand. It all started because of
      undocumented migrant workers who come into Thailand, and they are being
      taken advantage of because they are vulnerable and many fear being caught
      by the Thai police.}
  \item Who does the issue affect and how? \\
      \bf{The issue affects Thailand and other countries in Asia. The reason is
      because Thailand is bringing the people and selling them to other countries,
      which might end up making more unemployment because big corporations don't
      need any people to hire, they can just buy people.}
  \item How have the United States, other nations, and international
    organizations responded to the issue? \\
      \bf{In $2000$, the $106$th US Congress passed the \it{Trafficking Victims
      Protection Act (TVPA)}. It's the first comprehensive federal law that's
      designed to protect victims of sex and labor trafficking, to prosecute
      traffickers, and to prevent human trafficking in the US and around the
      world. That same year, the UN General Assembly adopted the UN Protocol to
      prevent trafficking of any kind and punish the human traffickers.}
  \item How do these responses relate to the foreign policy spectrum of
    approaches you studied in the lesson? (Recall that the approaches are
    isolationism, diplomacy, interventionism, imperialism. The response could
      relate to more than one approach.) \\
      \bf{In $2008$, Thailand passed the Anti-TIP act, which criminalizes human
      trafficking and provides any help to trafficking victims. Each MOU creates
      a task force in each country to work together to stop the trafficking
      between the countries.}
  \item Have the efforts and policies enacted by governments and international
    organizations been successful? Explain. \\
      \bf{No. Because when you look at the statistics, the numbers just keep
      increasing and they will keep increasing until somebody starts to take
      action.}
\end{itemize}
